\chapter{Virtual Memory}

Virtual memory treats main memory as a ``cache'' for secondary storage
(disk/SSD). Each program sees its own large, private \textbf{virtual address
space}; hardware and the operating system translate virtual addresses to
\textbf{physical addresses}.

\section{Why Virtual Memory?}

\begin{itemize}
  \item \textbf{Protection and isolation}: a process cannot read or corrupt
    another process's memory.
  \item \textbf{Sharing}: several programs share physical memory safely;
    shared libraries can be mapped into many address spaces.
  \item \textbf{Illusion of large memory}: programs can be larger than
    physical memory; unused parts stay on disk.
  \item \textbf{Relocation}: a program can be loaded anywhere in physical
    memory without changing its addresses.
\end{itemize}

\section{Pages and Address Translation}

Memory is divided into fixed-size \textbf{pages} (typically 4\,KB). The
virtual address is split into a \textbf{virtual page number (VPN)} and a
\textbf{page offset}; translation replaces the VPN by a \textbf{physical
page number (PPN)} and keeps the offset.

\begin{center}
\begin{tabular}{|c|c|}
\hline
Virtual page number & Page offset \\
\hline
\end{tabular}
\quad$\longrightarrow$\quad
\begin{tabular}{|c|c|}
\hline
Physical page number & Page offset \\
\hline
\end{tabular}
\end{center}

\begin{center}
\begin{tabular}{ll}
\toprule
Cache term & VM term \\
\midrule
block & page \\
miss & page fault \\
index/tag lookup & page table lookup \\
\bottomrule
\end{tabular}
\end{center}

\section{Page Tables}

\begin{definition}
A \textbf{page table} is an array in memory, indexed by VPN, whose entries
(PTEs) contain a \textbf{valid} bit, the PPN (or the disk location if not
resident), and status bits: \textbf{dirty}, \textbf{reference} (used),
and \textbf{protection} (read/write/execute, user/kernel). Each process
has its own page table; the hardware \textbf{page table register} points to
the current one.
\end{definition}

\begin{figure}[H]
\centering
\begin{tikzpicture}[font=\small, >=Stealth, thick, scale=0.9, transform shape]
  % virtual address
  \draw[fill=blue!10] (0,5) rectangle (4,5.7); \node at (2,5.35) {VPN (20)};
  \draw[fill=orange!15] (4,5) rectangle (6.4,5.7); \node at (5.2,5.35) {offset (12)};
  \node[anchor=east] at (-0.1,5.35) {virtual address};
  % page table
  \node[draw, fill=gray!15, font=\scriptsize, align=center] (ptr) at (-1.6,4.0) {page table\\register};
  \foreach \y in {0.8,1.3,...,4.3} \draw (1,\y) -- (4,\y);
  \draw (1,0.8) rectangle (4,4.3);
  \draw (1.5,0.8) -- (1.5,4.3);
  \node[font=\scriptsize] at (1.25,4.5) {V};
  \node[font=\scriptsize] at (2.75,4.5) {PPN};
  \fill[green!20] (1,2.3) rectangle (4,2.8);
  \draw (1,2.3) rectangle (4,2.8) (1.5,2.3) -- (1.5,2.8);
  \node[font=\scriptsize] at (1.25,2.55) {1};
  \node[font=\scriptsize] at (2.75,2.55) {PPN};
  \node[font=\scriptsize, below] at (2.5,0.8) {page table (in memory)};
  \draw[->] (ptr.east) -- (-0.2,4.0) |- node[pos=0.75, above, font=\scriptsize]{base} (1,4.3);
  \draw[->] (0.5,5) -- (0.5,2.55) -- (1,2.55);
  \node[font=\scriptsize, anchor=west] at (0.55,3.0) {index};
  % physical address
  \draw[fill=green!12] (5.5,1.3) rectangle (8.5,2); \node at (7,1.65) {PPN (18)};
  \draw[fill=orange!15] (8.5,1.3) rectangle (10.9,2); \node at (9.7,1.65) {offset (12)};
  \node[anchor=north] at (8.2,1.2) {physical address};
  \draw[->] (4,2.55) -- (4.75,2.55) -- (4.75,1.65) -- (5.5,1.65);
  \draw[->] (5.2,5) -- (5.2,4.2) -- (9.7,4.2) -- (9.7,2);
  \node[font=\scriptsize, above] at (8,4.2) {copied unchanged};
\end{tikzpicture}
\caption{Translation with a single-level page table (32-bit virtual
addresses, 4\,KB pages, a 30-bit physical address assumed for the picture).
The VPN indexes the table; the offset is not translated.}
\end{figure}

\begin{example}[Page table size]
32-bit virtual addresses, 4\,KB pages, 4-byte PTEs. Offset = 12 bits,
VPN = 20 bits, so $2^{20}$ entries $\times$ 4\,B = 4\,MB per process. For
64-bit addresses this is far too big, so real systems use
\textbf{multi-level (hierarchical) page tables}, where only the parts of the
table actually in use are allocated, or inverted/hashed page tables.
\end{example}

\begin{figure}[H]
\centering
\begin{tikzpicture}[font=\scriptsize, >=Stealth, thick]
  \draw[fill=blue!10] (0,4) rectangle (2,4.6); \node at (1,4.3) {P1 (10)};
  \draw[fill=blue!20] (2,4) rectangle (4,4.6); \node at (3,4.3) {P2 (10)};
  \draw[fill=orange!15] (4,4) rectangle (6,4.6); \node at (5,4.3) {offset (12)};
  % level 1
  \draw (0,0) rectangle (1.6,3.2); \node[below] at (0.8,0) {outer table};
  \foreach \y in {0.4,0.8,...,2.8} \draw (0,\y) -- (1.6,\y);
  \fill[green!20] (0,2.0) rectangle (1.6,2.4); \draw (0,2.0) rectangle (1.6,2.4);
  \node at (0.8,0.6) {invalid};
  \node at (0.8,1.0) {invalid};
  \draw[->] (1,4) -- (1,3.5) -- (-0.3,3.5) -- (-0.3,2.2) -- (0,2.2);
  % level 2
  \draw (3,0) rectangle (4.6,3.2); \node[below, align=center] at (3.8,0) {one second-level\\table};
  \foreach \y in {0.4,0.8,...,2.8} \draw (3,\y) -- (4.6,\y);
  \fill[green!20] (3,1.2) rectangle (4.6,1.6); \draw (3,1.2) rectangle (4.6,1.6);
  \node at (3.8,1.4) {PPN};
  \draw[->] (1.6,2.2) -- (2.3,2.2) -- (2.3,3.2) -- (3,3.2);
  \draw[->] (3,4) -- (3,3.5) -- (2.7,3.5) -- (2.7,1.4) -- (3,1.4);
  % physical page
  \draw[fill=green!8] (6,0.4) rectangle (7.6,2.4); \node[align=center] at (6.8,1.4) {physical\\page};
  \draw[->] (4.6,1.4) -- (6,1.4);
  \draw[->] (5,4) -- (5,3) -- (6.8,3) -- (6.8,2.4);
\end{tikzpicture}
\caption{A two-level page table splits the VPN in two. Outer entries for
unused regions are simply invalid, so their second-level tables are never
allocated.}
\end{figure}

\section{Page Faults}

If the PTE is invalid, a \textbf{page fault} exception occurs and the OS:
\begin{enumerate}
  \item finds the page on disk;
  \item chooses a physical page to replace (approximating LRU using
    reference bits, e.g.\ the clock algorithm); if the victim is dirty,
    writes it back to disk;
  \item reads the needed page into memory and updates the page table;
  \item restarts the faulting instruction.
\end{enumerate}

\begin{important}[Design consequences of huge miss penalties]
A page fault costs millions of cycles, so:
\begin{itemize}
  \item pages are large (amortize the disk access);
  \item placement is \textbf{fully associative} (any virtual page in any
    physical page) to minimize the miss rate;
  \item misses are handled in \emph{software} (clever algorithms are worth it);
  \item always \textbf{write-back}, never write-through.
\end{itemize}
\end{important}

\section{The TLB}

Without help, every memory access would need an extra memory access to read
the PTE, doubling the access time. The \textbf{translation lookaside
buffer (TLB)} is a small cache of recent translations (typically 16--512
entries, often fully associative). Each entry holds a tag (VPN), PPN, and
valid, dirty, reference and protection bits.

\begin{figure}[H]
\centering
\begin{tikzpicture}[thick, >=Stealth, font=\small, node distance=1.1cm and 1.6cm,
  box/.style={draw, rounded corners, minimum height=0.8cm, align=center}]
  \node[box] (va) {Virtual address};
  \node[box, right=of va] (tlb) {TLB};
  \node[box, below=of tlb] (pt) {Page table\\(memory)};
  \node[box, right=of tlb] (pa) {Physical address};
  \node[box, right=of pa] (c) {Cache /\\memory};
  \node[box, below=of pt] (os) {OS: page fault\\handler (disk)};
  \draw[->] (va) -- (tlb);
  \draw[->] (tlb) -- node[above, font=\scriptsize]{hit} (pa);
  \draw[->] (tlb) -- node[left, font=\scriptsize]{miss} (pt);
  \draw[->] (pt.east) -| node[pos=0.25, below, font=\scriptsize]{valid: refill TLB} (pa.south);
  \draw[->] (pt) -- node[left, font=\scriptsize]{invalid} (os);
  \draw[->] (pa) -- (c);
\end{tikzpicture}
\caption{Translation flow.}
\end{figure}

Possible outcomes of an access:
\begin{center}
\begin{tabular}{cccl}
\toprule
TLB & Page table & Cache & Possible? \\
\midrule
hit & hit & hit/miss & yes (best case: hit/hit/hit) \\
miss & hit & hit/miss & yes: refill TLB from page table \\
miss & miss & miss & yes: page fault \\
hit & miss & -- & no: TLB only caches valid translations \\
miss & miss & hit & no: data not in memory cannot be in the cache \\
\bottomrule
\end{tabular}
\end{center}

TLB misses can be handled in hardware (a page table walker, as in x86) or in
software (an exception handler, as in MIPS).

\section{Caches and Virtual Memory Together}

\begin{itemize}
  \item \textbf{Physically addressed cache}: translate first, then access the
    cache. Simple but the TLB is on the critical path.
  \item \textbf{Virtually indexed, physically tagged (VIPT)}: index the cache
    with page-offset bits (which do not change under translation) while the
    TLB translates in parallel; compare the physical tag afterwards. This
    works when index + offset bits fit in the page offset.
  \item Virtually addressed caches avoid translation but suffer from aliasing
    (two virtual addresses for the same physical data) and must be flushed
    or tagged with process IDs on context switch.
\end{itemize}

\begin{example}[Address breakdown]
32-bit virtual addresses, 30-bit physical addresses, 8\,KB pages. Offset =
13 bits; VPN = 19 bits; PPN = 17 bits. A 4-entry fully associative TLB stores
a 19-bit tag and a 17-bit PPN per entry (plus status bits).
\end{example}

\section{Protection}

The OS (running in kernel/supervisor mode) is the only one allowed to modify
page tables and the page table register. User programs enter the kernel via
\textbf{system calls} (\code{syscall} on MIPS), which switch modes. On a
context switch the OS changes the page table register and invalidates (or
tags with an address-space ID) the TLB.

\section*{Summary}
\begin{itemize}
  \item VM gives each process a private address space: protection, sharing,
    relocation, and the illusion of more memory.
  \item Virtual address = VPN | offset; page table maps VPN $\to$ PPN with
    valid, dirty, reference and protection bits.
  \item Page faults are very expensive: large pages, fully associative
    placement, LRU approximation, write-back, handled by the OS.
  \item The TLB caches translations; know which TLB/page table/cache
    combinations are possible.
\end{itemize}
